\section{人类语言与词义表示}

\subsection{语言：人类智能的核心驱动力}

Manning 教授从语言学家的视角出发，论证了语言在人类智能中的核心地位。他将人类与近亲黑猩猩进行了对比：

\begin{knowledgebox}{语言与人类文明}
黑猩猩在很多方面与人类相似——它们会使用工具、能规划解决问题、甚至短期记忆优于人类。但人类与黑猩猩之间最关键的区别就是\textbf{语言}。语言不仅是交流工具，还是高层次思维的支架（scaffolding），使人类能够进行更精细的推理和规划。
\end{knowledgebox}

语言在人类文明中扮演了三重角色：

\begin{enumerate}
    \item \textbf{交流工具}：使人类能够协作，实现``人类主宰''（human ascendancy）
    \item \textbf{思维支架}：为高层次认知提供结构化框架——我们用语言来思考、计划、推理
    \item \textbf{知识载体}：书写系统（约5000年历史）使知识能够跨越时间和空间传播，从青铜时代到现代科技
\end{enumerate}

\begin{figure}[H]
\centering
\includegraphics[width=0.95\textwidth]{slides-images/slide-009.jpg}
\caption{人类语言的多面性：作为社会工具的语言灵活而富有变化\protect\footnotemark}
\end{figure}
\footnotetext{Slides 第10页。}

Manning 还引用了 Stanford 心理学家 Herb Clark 的名言：

\begin{importantbox}{语言的本质}
``The common misconception is that language use has primarily to do with words and what they mean. It doesn't. It has primarily to do with people and what they mean.'' —— Herb Clark

语言使用的核心不在于词语及其含义，而在于人以及人想要表达的意思。
\end{importantbox}

\begin{warningbox}{语言不是静态的}
语言不是从天而降的固定系统——它由人类构造，并随每一代人不断变化。大部分语言创新发生在青少年和年轻人群体中。这意味着任何基于固定规则的 NLP 系统都面临过时的风险。
\end{warningbox}

\subsection{深度学习驱动的 NLP 革命}

过去十年，深度学习为 NLP 带来了革命性进步。Manning 列举了几个标志性成就：

\subsubsection{机器翻译}

神经机器翻译（Neural Machine Translation）从 2014 年开始发展，到 2016 年已被 Google 等服务大规模部署。这使得跨语言交流变得前所未有地容易。

\begin{figure}[H]
\centering
\includegraphics[width=0.95\textwidth]{slides-images/slide-011.jpg}
\caption{Google 翻译示例：将斯瓦希里语新闻翻译为英语，实现跨语言信息获取\protect\footnotemark}
\end{figure}
\footnotetext{Slides 第12页。}

\subsubsection{从搜索引擎到问答引擎}

传统搜索引擎基于关键词匹配返回文档列表。现代系统使用三层神经网络流水线实现真正的\textbf{问答}：

\begin{enumerate}
    \item \textbf{检索网络}（Retrieval）：找到与查询语义相似的段落
    \item \textbf{重排序网络}（Reranking）：对候选段落重新排序
    \item \textbf{阅读网络}（Reading）：从候选段落中综合信息并生成答案
\end{enumerate}

\subsubsection{GPT-2 与大语言模型的崛起}

2019 年的 GPT-2 首次展示了大语言模型生成流畅文本的能力。Manning 展示了一个经典示例：给定一个关于``辛辛那提被盗核材料''的开头，GPT-2 能生成连贯的新闻续写，不仅语法正确，还展现出丰富的世界知识——它知道辛辛那提在俄亥俄州、核材料由美国能源部监管等。

\begin{figure}[H]
\centering
\includegraphics[width=0.95\textwidth]{slides-images/slide-013.jpg}
\caption{GPT-2 文本生成示例：基于提示生成连贯的新闻故事，展现出语法正确性和世界知识\protect\footnotemark}
\end{figure}
\footnotetext{Slides 第14页。}

\begin{knowledgebox}{GPT-2 的工作原理}
GPT-2 的核心机制极其简单：给定所有已有文本，预测下一个最可能出现的词（next word prediction）。然后将预测的词加入已有文本，继续预测下一个词。这种简单的自回归生成方式，在足够大的模型和数据上，产生了惊人的效果。
\end{knowledgebox}

\subsubsection{ChatGPT 与指令遵循}

ChatGPT（及 GPT-4）的巨大突破在于\textbf{指令遵循}（instruction following）：用户可以用自然语言给出指令，模型能理解并执行。Manning 展示了让 ChatGPT 代写请假邮件的例子——即使输入中包含拼写错误（``song'' 而非 ``son''），模型也能正确理解意图。

\begin{figure}[H]
\centering
\includegraphics[width=0.95\textwidth]{slides-images/slide-014.jpg}
\caption{ChatGPT 指令遵循示例：正确理解含拼写错误的指令，生成得体的请假邮件\protect\footnotemark}
\end{figure}
\footnotetext{Slides 第15页。}

\subsubsection{多模态基础模型}

Stanford 提出的\textbf{基础模型}（Foundation Models）概念，将大语言模型的技术范式推广到图像、声音、生物信息学等多种模态。Manning 用 DALL-E 生成``火车通过金门大桥''的图片作为示例，展示了文本到图像生成的能力。

\begin{figure}[H]
\centering
\includegraphics[width=0.95\textwidth]{slides-images/slide-015.jpg}
\caption{DALL-E 3 生成的图像：根据详细的文本描述生成符合要求的场景插画\protect\footnotemark}
\end{figure}
\footnotetext{Slides 第16页。}

\subsection{本章小结}

人类语言是文明的核心驱动力，同时也是 AI 面临的最具挑战性的问题之一。深度学习在过去十年带来了 NLP 的革命性进步——从机器翻译到 GPT 系列大语言模型，再到多模态基础模型。这些进步的基础，正是本课程将要系统学习的技术。

%% ========== Part 3: 词义表示 ==========

